\uuid{5l6A}
\titre{Ensembles}
\auteur{bodin ; croizat ; sacre}
\organisation{exo7}
\licence{CC BY-NC-SA 4.0}
\datecreate{2022-03-08}
\niveau{L1}
\module{Algèbre}
\chapitre{Logique, ensemble, raisonnement}
\sousChapitre{Ensemble}
\difficulte{3}
\theme{Ensembles}
\exo7qcmid{hilisit.logique.039}

\contenu{
\begin{question}
\qclassification{100.02}
Soit la fonction réelle définie par $f(x) = x^2 - x - 2$. Pour quelle fonction $u$ a-t-on $f \circ u (x) = 9(x^2+x)$ ?
\begin{answers}
	\bad{$u(x) = 9x$}
	\good{$u(x) = 3x+2$}
	\bad{$u(x) = -3x$}
	\bad{$u(x) = 9x+2$}
\end{answers}
\begin{explanations}
On calcule pour $u(x) = 3x+2$ que $f \circ u (x) = f(3x+2) = (3x+2)^2 - (3x+2)-2 = 9x^2 + 12x + 4 - 3x - 2 - 2 = 9x^2 + 9x = 9(x^2+x)$.
\end{explanations}
\end{question}
}
